\BeleskeID{09_3-s019}
% element: 09_3-s019:e05
% slide-id: 09_3-s019
Dobijeni izraz može se posmatrati tako što prvo fiksiramo broj invertora. Za izabrano $N$ tražimo međusobne odnose dimenzija koji daju minimalno kašnjenje.

Prvi stepen ostaje minimalne geometrije i ima normiranu veličinu 1. Idealna veličina drugog stepena je $f$, trećeg $f^2$, a zatim se nastavlja geometrijski niz. Kada su potrebne celobrojne normirane dimenzije, veličine se zaokružuju. Za nenegativne vrednosti $\lfloor x+0.5\rfloor$ bira najbliži ceo broj.

Posle zaokruživanja treba računati kašnjenje iz stvarnih odnosa susednih veličina. Zaokružene veličine ne moraju dati potpuno jednake odnose kao idealan geometrijski niz.
